Conversely, let $y$ be a $W$-point of $\cI_T^{\mathrm{adm}}$, write
$p_W:C_T\times_T W\to W$ for the projection, put
$y_1=y$ and $y_2=-x_T-y$, and let $D_1,D_2$ be their relative Abel inverse
divisors.  Reading \eqref{eq:abel-class-sum} from right to left gives
\[
 \mathcal O(D_1+D_2)
 \simeq\mathcal O(3\kappa-D)\otimes p_W^*\mathcal N
\]
for a line bundle $\mathcal N$ on $W$.  The divisor section therefore gives
a unique projective section $[s]$ with
$\operatorname{div}(s)=D+D_1+D_2$; the factor pulled back from $W$ does not
affect this projective point.  Since $y_i\in P_T$,
\[
 \Nm_\pi\mathcal O(D_i)=\Nm_\pi(\kappa)=\eta^3.
\]
Thus the effective divisor $\pi_*D_i$ on $E_T$ is cut out by a unique
projective line $\ell_i$.  Multiplicativity of the finite locally free norm
gives, up to a
unit from $W$,
\[
 \Nm_\pi(s)=\ell_D\ell_1\ell_2,
 \qquad Q_s=\ell_1\ell_2,
\]
where $\ell_D$ is the fixed norm line of $D$.  Admissibility makes the two
residual factors distinct and orders them by the ordered pair $(D_1,D_2)$.
This constructs the inverse of \eqref{eq:abel-incidence-isomorphism}.

The relative Abel inverse on the $h^0=1$ locus, formation of zero divisors,
and finite locally free norms commute with arbitrary base change.  The two
composites preserve the universal Cartier divisors and hence are the identity
as morphisms.  Swapping $D_1,D_2$ gives $y_1\leftrightarrow y_2$, which is
the stated incidence involution.
\end{proof}

\begin{lemma}[Nonemptiness of the admissible locus]
\label[lemma]{lem:admissible-nonempty}
The open subscheme $\mathscr G$ of \cref{def:admissible-locus} is nonempty in
characteristic zero.  It is dense in $\mathscr R_T$, dominates $T$, and its
geometric generic fibre is a nonempty dense open of the geometric generic
residual curve.
\end{lemma}

\begin{proof}
For nonemptiness, use the following characteristic-$13$ point of the
normalized chart:
\[
 (c,m,e_1,e_2,e_3)=(1,1,1,0,2),\qquad \lambda=m^{-3}=1,
 \qquad
 (u:v:z)=(3:5:1).
\]
Direct substitution verifies all four conditions of
\cref{def:admissible-locus}: the fixed divisor is reduced; the three norm
lines cut transverse, pairwise disjoint sections away from the branch
divisor; each residual point lifts to one simple zero; and both residual
divisors lie in the smooth $h^0=1$ Abel locus.  Every coefficient,
discriminant, intersection point, and nonvanishing test in this substitution
is displayed in \cref{app:admissible-specialization}.

For the passage to characteristic zero, enlarge the localization
$R_0=\mathbf Z[1/N]$ from \cref{lem:residual-octic-integral}, without
inverting $13$, so that the chart and all admissibility conditions spread
out.  Work on the affine chart $v\ne0$.  The coordinate ring $A_0$ of a
suitable affine open in the spread of $T_{\QQ}$ is a localization of
\[
 R_0[c,c^{-1},m,m^{-1},e_1,e_2,e_3]
\]
and hence is a unique factorization domain.  With
$U_0=u/v$ and $Z_0=z/v$, clear denominators in $F_8(U_0,1,Z_0)$ and divide
by its content.  The resulting primitive polynomial
$f\in A_0[U_0,Z_0]$ is irreducible over $\operatorname{Frac}(A_0)$ by
\cref{lem:residual-octic-integral,lem:normalized-function-field}; Gauss's
lemma makes it irreducible in $A_0[U_0,Z_0]$.  Thus
\[
 \mathscr Y=\operatorname{Spec}A_0[U_0,Z_0]/(f)
\]
is integral.  Since $R_0$ injects into this domain, its coordinate ring is
torsion-free as an $R_0$-module; over the localization $R_0$ of $\mathbf Z$
this is equivalent to flatness.  Hence $\mathscr Y$ is $R_0$-flat.

Admissibility defines an open subscheme of $\mathscr Y$ containing the
displayed point over $\mathbf F_{13}$.  It therefore contains a principal
open $D(g)$ with $g\ne0$.  Since $\mathscr Y$ is integral, $D(g)$ contains
the generic point of $\mathscr Y$; in particular, its characteristic-zero
generic fibre is nonempty.  The generic residual curve dominates the generic
point of $T_{\QQ}$, and this remains true on $D(g)$.  After base change to
$\CC$, this open lies in $\mathscr G$.  Hence $\mathscr G$ is nonempty, dense in $\mathscr R_T$, and
dominates $T$; its geometric generic fibre is a nonempty open of a
geometrically integral curve and is therefore dense.
\end{proof}

\begin{corollary}[The two admissible branches]
\label[corollary]{cor:two-admissible-branches}
After replacing $\mathscr G$ by a dense open subset, there is an isomorphism
over $\mathscr G$
\[
 \cI_T^{\mathrm{adm}}\simeq\mathscr G\sqcup\mathscr G.
\]
The incidence involution exchanges the two copies.  On a geometric generic
fibre over $T$, each copy is a dense open of a geometrically integral curve.
\end{corollary}

\begin{proof}
By \cref{prop:ordering-cover}, the finite
\'etale factor-ordering cover has two rational sections over $\CC$ at the
generic point of $\mathscr G$.  They extend to two disjoint sections after
shrinking $\mathscr G$, so
$\operatorname{Ord}_{\mathrm{fac},\mathscr G}
 \simeq\mathscr G\sqcup\mathscr G$.
Apply \cref{prop:abel-incidence}.  The last assertion follows from
\cref{lem:residual-octic-integral,lem:admissible-nonempty}.
\end{proof}

\subsection{Boundary exhaustion and global labels}

Let $\xi$ be the generic point of the integral scheme $\cM$, and let
$\bar\xi$ be a geometric generic point.  For a geometric point $x$ of
$\cM$, define the pair fibre scheme-theoretically by
\[
 C_x=X\cap(-x-X)\subset P,
 \qquad C_x^{\mathrm{ne}}=C_x\setminus\{0,-x\}.
\]
Over $\bar\xi$, the unique effective divisors representing $x,y,-x-y$
satisfy
\begin{equation}
 D_x+D_y+D_{-x-y}=\operatorname{div}(s_y)
\label{eq:pair-zero-divisor}
\end{equation}
for a projective section $[s_y]\in\PP H^0(C,3\kappa-D_x)$.  This defines a
rational map
\[
 \Psi:C_x\dashrightarrow\PP H^0(C,3\kappa-D_x).
\]
\begin{lemma}[Boundary exhaustion]
\label[lemma]{lem:boundary-exhaustion}
 The rational map $\Psi$ has finite fibres wherever it is defined.  The
 determinant component $u=v$ gives no nonexceptional incidence-curve
 component; its generic divisor data give only the Abel classes $y=0,-x$.
 Every irreducible curve component of
$C_{\bar\xi}^{\mathrm{ne}}$ dominates the residual octic $\mathscr R$, meets
its admissible open, and is the closure of one of the two branches in
\cref{cor:two-admissible-branches}.  In particular,
$C_{\bar\xi}^{\mathrm{ne}}$ has exactly two irreducible components.
\end{lemma}
